\documentclass[11pt,a4paper]{article}
\usepackage[T1]{fontenc}
\usepackage{lmodern,amsmath,amssymb,amsthm,mathtools}
\usepackage[margin=25mm]{geometry}
\usepackage{microtype,booktabs,enumitem,xcolor,listings,xurl,hyperref,fancyhdr}
\hypersetup{colorlinks=true,linkcolor=blue!45!black,citecolor=blue!45!black,urlcolor=blue!45!black,pdftitle={Partial Results on Sun's A179524 Divisibility Conjecture}}
\setlength{\headheight}{14pt}
\pagestyle{fancy}
\fancyhf{}
\fancyhead[L]{Sun's A179524 divisibility conjecture}
\fancyhead[R]{Partial results}
\fancyfoot[C]{\thepage}
\newtheorem{theorem}{Theorem}[section]
\newtheorem{lemma}[theorem]{Lemma}
\newtheorem{corollary}[theorem]{Corollary}
\newtheorem{proposition}[theorem]{Proposition}
\theoremstyle{remark}\newtheorem{remark}[theorem]{Remark}
\newcommand{\Z}{\mathbb Z}
\newcommand{\F}{\mathbb F}
\newcommand{\Q}{\mathbb Q}
\newcommand{\CT}{\operatorname{CT}}
\newcommand{\rad}{\operatorname{rad}}
\newcommand{\vp}{v_p}
\lstset{language=Python,basicstyle=\ttfamily\footnotesize,keywordstyle=\color{blue!55!black},commentstyle=\color{gray!80!black},stringstyle=\color{gray!80!black},breaklines=true,columns=fullflexible,keepspaces=true,showstringspaces=false,frame=single,rulecolor=\color{gray!40},aboveskip=10pt,belowskip=10pt}
\title{\textbf{Partial Results on Sun's A179524\\Divisibility Conjecture}\\[5pt]\large A quadratic invariant, prime-square transfer, and complex multiplication}
\author{Research note prepared with ChatGPT for Keio}
\date{7 November 2026}
\begin{document}
\maketitle
\begin{abstract}
Let \(a_n=\sum_k(-4)^k\binom nk^2\binom{n-k}k^2\) and
\(S(n)=\sum_{j=0}^{n-1}(20j+17)a_j\). Sun conjectured in 2010 that
\(n\mid S(n)\) for every positive integer \(n\). We do not prove that full
weighted assertion. We do prove the complete inert unweighted branch,
namely that the sum of the first $p$ terms is zero modulo $p^2$ for
$p\equiv11,13,17,19\pmod{20}$, using a quadratic differential invariant.
We also prove \(\rad(n)\mid S(n)\) and the exact identity
\(v_2(S(n))=v_2(n)\) for every \(n\). A first-order digit expansion gives
\(S(p^2m)\equiv U(m)S(p^2)\pmod {p^2}\), where
\(U(m)=\sum_{j<m}a_j\), for every odd prime \(p\) and every positive \(m\).
Using a CM elliptic curve of discriminant \(-20\), we prove
\(p^2\mid S(p^2m)\) whenever \(p\equiv11,13,17,19\pmod {20}\).
We also give a symbolic proof for \(p=3\). Exact computation supplies the
same all-cofactor conclusion for every odd prime \(p\le67\).
Consequently the full conjecture holds on an explicit set of integers
of natural density at least \(0.9481679114\). This density measures a
proved domain, not the proportion of the mathematical difficulty resolved.
The remaining higher odd prime powers, experimental strengthening, and
limits of the priority investigation are stated explicitly.
\end{abstract}

\section{The unchanged question and the proved domain}

Use the convention that a binomial coefficient with an inadmissible lower
index is zero. Define
\begin{equation}\label{def}
 \begin{aligned}
 M(n,k)&=\binom nk\binom{n-k}k,\qquad
 a_n=\sum_{0\le2k\le n}(-4)^k M(n,k)^2,\\
 S(n)&=\sum_{j<n}(20j+17)a_j .
 \end{aligned}
\end{equation}
For \(0\le2k\le n\), the factorial expression is
\(M(n,k)=n!/[k!^2(n-2k)!]\).
The initial values are \(1,1,-15,-143,1,12801,100401,-555855\).
The unchanged assertion is
\begin{equation}\label{target}
                     n\mid S(n)\qquad(n\ge1).
\end{equation}
This is Conjecture 5.8(ii), printed p.~33, of Sun's 2010 preprint
\cite{Sun2010}, with \(a_n=W_n(1,-4)\). The OEIS record
A179524 attributes it to 1 July 2010 \cite{OEIS}. The latest version of
arXiv:1006.2776 was withdrawn to merge material into other papers; the
numbered source cited here is explicitly version 11.

For an integer \(N\ne0\), \(v_p(N)\) denotes its \(p\)-adic valuation.
Let \(\mathcal P_0\) be the eighteen odd primes at most \(67\), and put
\[
 \mathcal P_*=\mathcal P_0\cup
 \{p\text{ prime}:p\equiv11,13,17,19\pmod {20}\}.
\]
Define \(e_p=2\) for \(p\in\mathcal P_*\) and \(e_p=1\) for other odd primes.

\begin{theorem}[Combined partial result]\label{main}
For every positive integer \(n\),
\begin{equation}\label{divisor}
 v_2(S(n))=v_2(n),\qquad
 2^{v_2(n)}\prod_{\substack{p\mid n\\p\ {\rm odd}}}
 p^{\min(v_p(n),e_p)}\ \mid\ S(n).
\end{equation}
In particular, \eqref{target} holds if \(v_p(n)\le e_p\) for every odd
prime \(p\). There is no restriction on the exponent of \(2\).
\end{theorem}

The prime cases in this theorem use Sun's established Ramanujan-type
supercongruence \cite{SunRefinement}. The inert-prime-square cases use
standard CM reduction and the Hasse invariant. The finite set
\(\mathcal P_0\) uses the exact certificate in Appendix~A, together with
an all-cofactor transfer theorem. No computational bound is imposed on
the cofactor \(m\) in any of these conclusions.

\begin{corollary}[An explicit density lower bound]\label{density}
The conjecture holds on a set of natural density at least
\[
 \frac8{\pi^2}\prod_{\substack{3\le p\le67\\p\ {\rm prime}}}
 \frac{1-p^{-3}}{1-p^{-2}}
 =
 \frac8{\pi^2}\prod_{\substack{3\le p\le67\\p\ {\rm prime}}}
 \frac{p^2+p+1}{p(p+1)}
 =0.94816791146370398749\ldots .
\]
\end{corollary}
\begin{proof}
Take the smaller set defined by \(v_p(n)\le2\) for odd \(p\le67\)
and \(v_p(n)\le1\) for \(p>67\), with arbitrary \(v_2(n)\).
The Euler product for these valuation restrictions is the displayed
expression. Its existence follows by first imposing finitely many local
conditions, using the Chinese remainder theorem, and bounding the omitted
tail by \(\sum_{p>y}p^{-2}\). The additional inert-prime cases can only
enlarge the proved set.
\end{proof}

\section{Prime divisibility and Lucas propagation}

We first record two integer identities:
\begin{align}
 \binom jm^2
 &=\sum_{i=0}^m\binom mi\binom{m+i}i\binom j{m+i},
 \label{product}\\
 \sum_{j<n}(20j+17)\binom jt
 &=(20t+17)\binom n{t+1}
       +20(t+1)\binom n{t+2}. \label{hockey}
\end{align}
The first counts ordered pairs of \(m\)-subsets by their union size.
The second follows from the hockey-stick identity and
\(j\binom jt=t\binom jt+(t+1)\binom j{t+1}\).

\begin{lemma}[Ramanujan reduction]\label{primelemma}
For an odd prime \(p\), \(h=(p-1)/2\), and \(d=\lfloor p/4\rfloor\),
\begin{equation}\label{ramred}
 S(p)\equiv-(-1)^h\sum_{j=0}^{d}
       (20j+3)\frac{(4j)!}{j!^4(-1024)^j}\pmod p .
\end{equation}
\end{lemma}
\begin{proof}
For \(0\le k\le h\), the polynomial \(F_k(X)=\binom X{2k}^2\)
is defined over \(\F_p\) and satisfies
\(F_k(2k-1-X)=F_k(X)\). Thus
\[
 \sum_{r\in\F_p}(20r+17)F_k(r)
 =(20k+7)\sum_{r\in\F_p}F_k(r).
\]
Apply \eqref{product} and sum over \(0\le r<p\). Modulo \(p\), the
only surviving binomial \(\binom p{2k+i+1}\) has \(2k+i=p-1\).
Since \(i=p-1-2k\) is even, its coefficient reduces to
\(\binom{2k}{p-1-2k}\). Therefore
\[
 S(p)\equiv\sum_{k=\lceil h/2\rceil}^{h}
 (20k+7)(-4)^k\binom{2k}k^2\binom{2k}{2h-2k}\pmod p .
\]
Put \(k=h-j\). The needed congruences are
\[
 20k+7\equiv-(20j+3),\quad
 (-4)^{h-j}\equiv(-1)^h(-4)^{-j},
\]
\[
 \binom{p-1-2j}{h-j}^{\!2}
 \equiv\binom{h+j}{2j}^{\!2}
 \equiv\frac{\binom{2j}j^2}{256^j},\qquad
 \binom{p-1-2j}{2j}\equiv\binom{4j}{2j}.
\]
For the middle identity use
\(\binom{j-1/2}{2j}=(-1)^j\binom{2j}j/16^j\).
All denominators are units modulo \(p\). These formulas give
\eqref{ramred}.
\end{proof}

Sun's Theorem 1.2 \cite{SunRefinement} states, for \(p>3\),
\[
 \sum_{j=0}^{h}(20j+3)\frac{(4j)!}{j!^4(-1024)^j}
 \equiv(-1)^h p\bigl(2^{p-1}+2-(2^{p-1}-1)^2\bigr)\pmod {p^4}.
\]
For \(d<j\le h\), the numerator \((4j)!\) is divisible by \(p\)
and \(j!\) is not. Hence Lemma~\ref{primelemma} gives \(p\mid S(p)\).
The remaining primes satisfy \(S(2)=54\) and \(S(3)=-801\).

\begin{lemma}[Digit congruence]\label{lucas}
For every odd prime \(p\), \(u\ge0\), and \(0\le r<p\),
\[
                 a_{pu+r}\equiv a_u a_r\pmod p.
\]
\end{lemma}
\begin{proof}
Write the inner index as \(k=pv+s\), \(0\le s<p\).
Multinomial Lucas makes \(M(pu+r,pv+s)\) vanish modulo \(p\)
unless \(2s\le r\). In the latter case
\[
 M(pu+r,pv+s)\equiv M(u,v)M(r,s)\pmod p.
\]
The high-digit range is exactly \(0\le2v\le u\).
Also \((-4)^{pv+s}\equiv(-4)^{v+s}\pmod p\).
Summing the factored terms proves the claim.
\end{proof}

For every \(m\ge1\), it follows that
\begin{equation}\label{primeprop}
 S(pm)\equiv S(p)\sum_{u<m}a_u\equiv0\pmod p.
\end{equation}
The case \(p=2\) follows directly from \(a_u\equiv1\pmod2\).
Consequently \(\rad(n)\mid S(n)\) for every \(n\).
This argument alone does not increase the prime-power exponent.

\section{The entire 2-part, with exact valuation}

\begin{theorem}\label{twoadic}
For every \(n\ge1\), \(S(n)/n\) is an odd \(2\)-adic integer.
Thus \(S(n)\ne0\) and \(v_2(S(n))=v_2(n)\).
\end{theorem}
\begin{proof}
Rewrite \(M(j,k)=\binom j{2k}\binom{2k}k\).
For a fixed \(k\ge1\), use \eqref{product} with \(m=2k\), followed
by \eqref{hockey}, and divide by \(n\). Its contribution is an integer
linear combination of
\[
 (-4)^k\binom{2k}k^2
 \left(
 \frac{20t+17}{t+1}\binom{n-1}t+
 \frac{20(t+1)}{t+2}\binom{n-1}{t+1}
 \right),\qquad 2k\le t\le4k.
\]
The prefactor has \(2\)-adic valuation at least \(2k+2\), because
\(\binom{2k}k=2\binom{2k-1}{k-1}\).
Both displayed denominators are at most \(4k+2<2^{2k+2}\).
Every such term therefore has valuation at least one.
All sums are finite. The \(k=0\) contribution is
\[
                  \frac1n\sum_{j<n}(20j+17)=10n+7,
\]
which is odd. This proves the assertion, including \(n=1\).
\end{proof}

\section{A prime-square transfer valid for every cofactor}

Put \(f(n,k)=(-4)^kM(n,k)^2\), and define
\[
 L_n=\sum_k kf(n,k),\qquad
 U(N)=\sum_{n<N}a_n,\quad T(N)=\sum_{n<N}n a_n,\quad
 V(N)=\sum_{n<N}L_n.
\]
Thus \(S(N)=20T(N)+17U(N)\).
For an odd prime \(p\), let
\[
 \begin{aligned}
 Q_p&=\frac{(-4)^{p-1}-1}{p},\\
 A_r&=\sum_s f(r,s)(H_r-H_{r-2s}),\qquad
 B_r=\sum_s f(r,s)(H_{r-2s}-H_s)\quad(0\le r<p),
 \end{aligned}
\]
where \(H_0=0\), \(H_j=\sum_{i=1}^j1/i\), and only admissible \(s\)
are included. These harmonic numbers are \(p\)-integral.

\begin{lemma}[First-order digit expansion]\label{firstorder}
For every \(q\ge0\) and \(0\le r<p\),
\begin{equation}\label{digit2}
 a_{pq+r}\equiv a_q a_r+
 p\{2q a_q A_r+L_q(Q_p a_r+4B_r)\}\pmod {p^2}.
\end{equation}
\end{lemma}
\begin{proof}
Let \(E_p(N)=\prod_{1\le i\le N,\ p\nmid i}i\).
The exact factorization
\(N!=p^{\lfloor N/p\rfloor}\lfloor N/p\rfloor!E_p(N)\), together with
\(\sum_{i=1}^{p-1}i^{-1}=0\) in \(\F_p\), gives
\[
 E_p(pq+r)\equiv((p-1)!)^q r!(1+pqH_r)\pmod {p^2}.
\]
Write \(k=pj+s\). If \(2s>r\), Lucas gives \(p\mid M(pq+r,k)\),
so its square vanishes modulo \(p^2\).
Otherwise the factorial factorization gives
\[
 M(pq+r,pj+s)\equiv M(q,j)M(r,s)
 [1+p\{qH_r-2jH_s-(q-2j)H_{r-2s}\}]\pmod {p^2}.
\]
Only \(E_p\)-factors and small factorials are inverted.
The integer \(M(q,j)\) is retained exactly, so higher carries cause
no problem. Finally
\[
 (-4)^{pj+s}\equiv(-4)^{j+s}(1+pjQ_p)\pmod {p^2}.
\]
Square the multinomial congruence, multiply, and sum over \(j,s\).
\end{proof}

\begin{theorem}[Prime-square transfer]\label{transfer}
For every odd prime \(p\) and every positive integer \(m\),
\begin{equation}\label{transferformula}
                    S(p^2m)\equiv U(m)S(p^2)\pmod {p^2}.
\end{equation}
\end{theorem}
\begin{proof}
Write \(w_r=20r+17\), \(U_p=U(p)\), \(T_p=T(p)\), and \(V_p=V(p)\).
Summing \eqref{digit2} with weight \(20(pq+r)+17\) gives
\begin{equation}\label{block}
 S(pM)\equiv S(p)U(M)+pC_pT(M)+pD_pV(M)\pmod {p^2},
\end{equation}
where
\[
 C_p=20U_p+2\sum_{r<p}w_rA_r,\qquad
 D_p=4\sum_{r<p}w_rB_r .
\]
The omitted term is exactly \(pQ_p S(p)V(M)\); it vanishes because
\(p\mid S(p)\), proved in Section 2.
Lucas also gives \(L_{pu+r}\equiv a_uL_r\pmod p\), since the index
\(pv+s\) reduces to \(s\). Hence
\[
 U(pm)\equiv U_pU(m),\quad T(pm)\equiv T_pU(m),\quad
 V(pm)\equiv V_pU(m)\pmod p.
\]
Substitute \(M=pm\) into \eqref{block}. Its three coefficients
\(S(p),pC_p,pD_p\) are divisible by \(p\), so these substitutions
are valid modulo \(p^2\). The result is
\[
 S(p^2m)\equiv U(m)[S(p)U_p+pC_pT_p+pD_pV_p].
\]
The bracket is \(S(p^2)\) modulo \(p^2\) by \eqref{block} with \(M=p\).
\end{proof}

\begin{remark}[What a finite base verifies]
For a fixed prime \(p\), an exact verification of \(p^2\mid S(p^2)\)
now proves \(p^2\mid S(p^2m)\) for \emph{every} \(m\).
It is a proof for an unbounded cofactor, but not a verification over
all primes. The theorem also gives no \(p^3\) conclusion.
\end{remark}

\section{An infinite prime-square family from complex multiplication}

\begin{theorem}\label{inert}
If \(p\equiv11,13,17,19\pmod {20}\) is prime, then
\[
                     p^2\mid S(p^2m)\qquad(m\ge1).
\]
\end{theorem}
\begin{proof}
For an indeterminate \(c\), define
\[
 W_r(c)=\sum_k c^kM(r,k)^2,\qquad
 \mathcal U_p(c)=\sum_{r<p}W_r(c).
\]
Set \(h=(p-1)/2\) and \(d=\lfloor p/4\rfloor\).
The binomial-square summation in Section 2, followed by \(k=h-j\),
gives the polynomial identity
\begin{equation}\label{Upoly}
 \mathcal U_p(c)\equiv
 c^h\sum_{j=0}^{d}\frac{(4j)!}{j!^4(256c)^j}\pmod p.
\end{equation}
All apparent negative powers cancel. The finite Legendre identity
\[
 P_d(t)^2=\sum_{j=0}^d\binom dj\binom{d+j}j\binom{2j}j
                   \left(\frac{t^2-1}{4}\right)^j
\]
and \(\{-d,d+1\}=\{1/4,3/4\}\) in \(\F_p\) imply
\begin{equation}\label{square}
 \mathcal U_p(c)\equiv c^hP_d(t)^2,\qquad t^2=1-\frac1c.
\end{equation}
This is an identity in the parameter, not merely an equality at a
particular numerical value.

Take \(t_0=\sqrt5/2\), corresponding to \(c=-4\), and consider
\[
 E_{t_0}:\quad y^2=x\left(x^2+2x+\frac{1-t_0}{2}\right).
\]
Its invariants are
\[
 \Delta=8(t_0-1)^2(t_0+1),\qquad
 j(E_{t_0})=
 \frac{64(3t_0+5)^3}{(t_0-1)^2(t_0+1)}
 =632000+282880\sqrt5.
\]
Duke \cite{Duke} identifies this value as \(j(i\sqrt5)\); hence the
curve has complex multiplication of discriminant \(-20\).
The other sign of \(t_0\) gives the conjugate CM value.
The coefficients are integral away from \(2\), and
\((t_0-1)(t_0+1)=1/4\), so the model has good reduction at every
prime above an odd rational prime.

The stated residue classes are exactly
\(\left(\frac{-5}{p}\right)=-1\). Deuring's reduction theorem therefore
makes this reduction supersingular; see \cite{EOY}.
This holds over \(\overline{\F}_p\), including when \(t_0\notin\F_p\).

For completeness, the Hasse coefficient can be read off directly.
Let \(b=(1-t)/2\). The coefficient of \(x^{p-1}\) in
\([x(x^2+2x+b)]^h\) is
\[
 \sum_{k=0}^{d}\frac{h!}{k!^2(h-2k)!}\,2^{h-2k}b^k
 =
 2^h\sum_{k=0}^{d}
 \frac{(-h/2)_k((1-h)/2)_k}{k!^2}b^k
 \equiv2^hP_d(t)\pmod p.
\]
Here \((z)_k\) is the rising factorial.
The standard Hasse-invariant criterion \cite[Theorem V.4.1(a)]{Silverman}
therefore gives \(P_d(t_0)=0\). Equation \eqref{square} yields
\begin{equation}\label{Uzero} U_p=\mathcal U_p(-4)\equiv0\pmod p.\end{equation}

There are two elementary ways to obtain the other moments.
The reflection used in Section 2 also gives, for each \(k\),
\[
 \sum_{r<p}r\binom r{2k}^2
 \equiv\left(k-\frac12\right)\sum_{r<p}\binom r{2k}^2\pmod p.
\]
Thus \(T_p\equiv V_p-U_p/2\pmod p\).
Since \(20T_p+17U_p=S(p)\equiv0\pmod p\) and \(p\ne2,5\),
\eqref{Uzero} implies \(T_p=V_p=0\) in \(\F_p\).
Alternatively, differentiate \eqref{square}: the square factor gives
\(\mathcal U_p'(-4)=0\) because \(t_0\ne0\); then
\(V_p=-4\mathcal U_p'(-4)=0\).

Finally the bracket in the proof of Theorem~\ref{transfer} vanishes
modulo \(p^2\): \(p\mid S(p)\) and \(U_p,T_p,V_p\) all vanish modulo \(p\).
This proves the theorem for every cofactor \(m\).
\end{proof}

\section{A complete unweighted branch: a quadratic invariant}\label{zeroliftsec}

The following statement is separate from the weighted all-\(n\) target.
It supplies a full proof of the zero branch of Sun's Conjecture 5.8(i)
\cite{Sun2010}, recorded as \texttt{OeisA179524.conjecture3} in
\cite{Formal}. Its historical priority is not asserted.

\begin{theorem}[Automatic lift away from singular parameters]\label{zerolift}
Let \(p\) be an odd prime and set
\(Y_p(c)=\sum_{r=0}^{p-1}W_r(c)\). If \(c_0\in\Z_p\),
\(c_0(c_0-1)\in\Z_p^\times\), and \(Y_p(c_0)\equiv0\pmod p\), then
\[
                         Y_p(c_0)\equiv0\pmod {p^2}.
\]
\end{theorem}

The binomial reductions used below are established ingredients:
Sun \cite[Theorem 1.1]{SunApery} relates this sum to a sum of Ap\'ery
polynomials, and Guo--Zeng \cite[Theorem 1.3]{GuoZeng} give the
binomial reduction used here. We include its direct proof to make the
subsequent differential argument checkable without additional files.
The finite-invariant lifting strategy itself has a direct antecedent in
Babanskyy \cite[Lemmas 2.1--2.2]{Babanskyy}: a quadratic differential
invariant, explicit treatment of exceptional degree-\(p\) coefficients,
and a double-root argument yield a lift from modulo \(p\) to modulo
\(p^2\). Here that strategy is adapted to the dual polynomial below.
The proof is given directly; the invariant principle is not claimed
as a new invention.

\begin{lemma}\label{vandermonde}
For an even integer \(m\) with \(0\le m\le p-1\),
\[
 \sum_{r=m}^{p-1}\binom rm^2
                \equiv\binom{p+m}{2m+1}\pmod {p^2}.
\]
\end{lemma}
\begin{proof}
For \(m\le r<p\), all denominators in \(H_r-H_{r-m}\) are \(p\)-adic
units, and
\[
 \binom{p+m-1-r}{m}\equiv
 \binom rm[1-p(H_r-H_{r-m})]\pmod {p^2}.
\]
The sign is positive because \(m\) is even.
The involution \(r\mapsto p+m-1-r\) permutes \(m,\ldots,p-1\),
preserves \(\binom rm^2\) modulo \(p\), and reverses
\(H_r-H_{r-m}\) modulo \(p\). The last fact follows from
\(H_{p-1-u}\equiv H_u\pmod p\). Thus the sum of the first-order
corrections is zero modulo \(p\).
Vandermonde's identity gives
\[
 \sum_{r=m}^{p-1}\binom rm\binom{p+m-1-r}{m}
                  =\binom{p+m}{2m+1},
\]
proving the claim, including \(m=0\) and \(m=p-1\).
\end{proof}

\begin{lemma}[Dual polynomial and differential equation]\label{dual}
Put \(h=(p-1)/2\) and
\[
 Z_p(c)=\sum_{s=0}^{h}v_s c^s,\qquad
 v_s=p\frac{(2s)!^4}{s!^4(4s+1)!}.
\]
Then \(v_s\in\Z_p\), \(Y_p\equiv Z_p\pmod {p^2}\) coefficientwise,
and, with \(\theta=c\,d/dc\),
\begin{equation}\label{ode}
 \mathcal L Y_p\equiv0\pmod {p^2},\qquad
 \mathcal L=\theta(4\theta-1)(4\theta+1)-2c(2\theta+1)^3 .
\end{equation}
\end{lemma}
\begin{proof}
The denominator \((4s+1)!\) contains at most one factor \(p\),
since \(4s+1\le2p-1\); it is canceled by the displayed factor \(p\).
Lemma~\ref{vandermonde} makes the coefficient of \(c^s\) in \(Y_p\)
congruent to
\[
 \binom{2s}s^2\binom{p+2s}{4s+1}
 =
 p\binom{2s}s^2
       \frac{\prod_{j=1}^{2s}(p^2-j^2)}{(4s+1)!}.
\]
The product is \((2s)!^2\) modulo \(p^2\), since \(2s\) is even.
Before division the error has valuation at least three; the factorial
denominator loses at most one. Hence the error is still zero modulo
\(p^2\), proving the coefficientwise congruence.

The exact coefficient relation is
\[
 s(4s-1)(4s+1)v_s=2(2s-1)^3v_{s-1}\qquad(1\le s\le h).
\]
Consequently the entire truncation residual is
\[
             \mathcal L Z_p=-2p^3v_h c^{h+1}.
\]
It vanishes modulo \(p^3\). Applying the integer differential operator
to \(Y_p-Z_p\) proves \eqref{ode}. No coefficient \(s\) or \(4s\pm1\)
has been inverted modulo \(p\).
\end{proof}

\begin{lemma}[Quadratic differential invariant]\label{invariant}
For a polynomial \(Y\), define
\[
 \mathcal J[Y]=16c^2(c-1)(Y'^2-2YY'')
              -16c(3c-2)YY'-(4c+1)Y^2.
\]
Then
\[
             \frac{d}{dc}\mathcal J[Y]=2Y\,\frac{\mathcal LY}{c}.
\]
In particular, \(\mathcal J[Y_p]\equiv0\pmod {p^2}\) coefficientwise.
\end{lemma}
\begin{proof}
The differential identity follows by expanding
\[
 \frac{\mathcal LY}{c}
 =16c^2(1-c)Y'''+(48c-72c^2)Y''
                  +(15-52c)Y'-2Y.
\]
The right side is polynomial. Thus \(\mathcal J[Y_p]'\) is zero
modulo \(p^2\). A zero derivative alone would not suffice, because a
degree-\(p\) term could survive. We check that term explicitly.

Since \(\deg Y_p=h\), the degree of \(\mathcal J[Y_p]\) is at most
\(2h+1=p\). Its constant coefficient is \(-Y_p(0)^2=-p^2\).
If \(K\) is the leading coefficient of \(Y_p\), its degree-\(p\)
coefficient is
\[
 [16(h^2-2h(h-1))-48h-4]K^2
                   =-4(2h+1)^2K^2=-4p^2K^2.
\]
Both endpoints vanish modulo \(p^2\). For every intermediate index
\(1\le j\le p-1\), the derivative gives \(j\, [c^j]\mathcal J[Y_p]=0\)
modulo \(p^2\), and \(j\) is a unit. All coefficients therefore vanish.
\end{proof}

\begin{proof}[Proof of Theorem~\ref{zerolift}]
Equation \eqref{square} gives
\[
 Y_p(c)\equiv c^hP_d(t)^2\pmod p,\qquad
 d=\lfloor p/4\rfloor,\quad t^2=1-1/c.
\]
Choose \(t_0\) over the algebraic closure of the residue field with
\(t_0^2=1-1/c_0\).
The assumptions make \(t_0\ne0\), \(t_0^2\ne1\), and \(P_d(t_0)=0\).
Such a Legendre root is simple. Indeed, if its multiplicity were
\(m\ge2\), the lowest-degree term of
\[
        (1-t^2)P_d''-2tP_d'+d(d+1)P_d=0
\]
would have a nonzero coefficient proportional to
\((1-t_0^2)m(m-1)\), since \(m\le d<p\).
The degree-zero case has no roots and is vacuous.

Because \(dt/dc=1/(2tc^2)\), differentiating the polynomial congruence
gives
\[
 Y_p'(c_0)\equiv0,\qquad
 Y_p''(c_0)\equiv
 2c_0^hP_d'(t_0)^2(2t_0c_0^2)^{-2}\ne0\pmod p.
\]
The expression is in the base residue field, as is the left side;
the field extension is used only to describe the root.

Evaluate the invariant at \(c_0\). Both \(Y_p(c_0)\) and
\(Y_p'(c_0)\) are divisible by \(p\), so every term except the
\(Y_pY_p''\) term vanishes modulo \(p^2\). We obtain
\[
       32c_0^2(c_0-1)Y_p(c_0)Y_p''(c_0)\equiv0\pmod {p^2}.
\]
Its coefficient multiplying \(Y_p(c_0)\) is a \(p\)-adic unit.
This forces \(Y_p(c_0)\equiv0\pmod {p^2}\).
\end{proof}

\begin{corollary}[Sun's inert unweighted branch]\label{unweighted}
For every prime \(p\equiv11,13,17,19\pmod {20}\),
\[
                 \sum_{r=0}^{p-1}a_r\equiv0\pmod {p^2}.
\]
\end{corollary}
\begin{proof}
The CM argument in Theorem~\ref{inert} gives \(Y_p(-4)=0\pmod p\).
Here \(c_0=-4\) satisfies \(c_0(c_0-1)=20\), a unit at all the
specified primes. Apply Theorem~\ref{zerolift}.
\end{proof}

The corollary closes all quantifiers of this particular unweighted
branch. It neither proves the split-prime evaluation branches of
Sun's Conjecture 5.8(i), nor the weighted all-\(n\) statement
\eqref{target} in Conjecture 5.8(ii).


\section{A direct symbolic proof modulo 9}

\begin{proposition}\label{nine}
For every \(m\ge1\), \(9\mid S(9m)\).
\end{proposition}
\begin{proof}
The digit expansion \eqref{digit2} at \(p=3\), or a direct factorial
expansion, gives for every \(q\ge0\)
\begin{equation}\label{ternary}
 a_{3q}\equiv a_q+6L_q,\qquad
 a_{3q+1}\equiv a_q,\qquad
 a_{3q+2}\equiv-15a_q+3L_q\pmod9.
\end{equation}
Here is an explicit check of the coefficients. The only low-digit
pairs with \(2s\le r\) and their corresponding weights are
\[
 (r,s)=(0,0),(1,0),(2,0),(2,1),\qquad f(r,s)=1,1,1,-16.
\]
Use \(H_0=0,H_1=1,H_2=3/2\), together with
\[
 (-4)^{3j+s}\equiv(-4)^{j+s}(1+6j)\pmod9.
\]
The correction divided by \(3\) is
\[
 2q(H_r-H_{r-2s})+4j(H_{r-2s}-H_s)+2j .
\]
Its \(q\)-coefficient is zero modulo \(3\) in all four cases, and its
\(j\)-coefficients are \(2,0,2,1\), respectively.
All low-carry terms vanish modulo \(9\); higher carries remain in
\(M(q,j)\). This proves \eqref{ternary} for arbitrary \(q\).

Consequently
\[
 \sum_{r=0}^2(20(3q+r)+17)a_{3q+r}
 \equiv3(q a_q+L_q)\pmod9.
\]
Modulo \(3\), Lucas gives \(a_{3u+r}\equiv a_u a_r\) and
\(L_{3u+r}\equiv a_u L_r\). Since
\[
 (a_0,a_1,a_2)\equiv(1,1,0),\qquad
 (L_0,L_1,L_2)\equiv(0,0,2)\pmod3,
\]
we have
\[
 \sum_{r=0}^2((3u+r)a_{3u+r}+L_{3u+r})
 \equiv a_u\sum_{r=0}^2(r a_r+L_r)=0\pmod3.
\]
Sum the preceding modulo-\(9\) block over \(q<3m\), and group the
right side into these three-term blocks.
\end{proof}

\section{Exact computation and an equivalent unsolved condition}

Appendix A evaluates the original sums using integer arithmetic.
For fixed \(n\), its summands satisfy
\begin{equation}\label{termrec}
 f(n,0)=1,\qquad
 f(n,k+1)=
 \frac{-4(n-2k)^2(n-2k-1)^2}{(k+1)^4}f(n,k).
\end{equation}
Every division is checked to be exact. The run verified
\(n\mid S(n)\) for \(1\le n\le5000\), with no counterexample.
In particular \(p^2\mid S(p^2)\) was verified for
\[
 \mathcal P_0=\{3,5,7,11,13,17,19,23,29,31,37,41,43,47,53,59,61,67\}.
\]
Theorem~\ref{transfer} converts these eighteen exact base checks into
\(p^2\mid S(p^2m)\) for all \(m\), proving the finite-prime part of
Theorem~\ref{main}. This completes that theorem and its density corollary.

Define the Newton coefficients
\[
 b_\ell=\sum_{j=0}^{\ell}(-1)^{\ell-j}\binom{\ell}j a_j,\qquad b_{-1}=0.
\]
Binomial inversion and \eqref{hockey} give the exact formula
\begin{equation}\label{newton}
 \frac{S(n)}n=\sum_{\ell<n}\binom{n-1}{\ell}
 \frac{(20\ell+17)b_\ell+20\ell b_{\ell-1}}{\ell+1}.
\end{equation}
Since the binomial transform and its inverse have integer coefficients,
the full conjecture is equivalent to
\begin{equation}\label{equivalent}
                \ell+1\mid3b_\ell+20b_{\ell-1}\quad(\ell\ge0).
\end{equation}
This is a reformulation, not a proof of the missing cases.
A useful constant-term model is
\[
 P=(1+x+y)\left(1+\frac zx-\frac4{yz}\right),\quad Q=P-1,
 \qquad a_n=\CT P^n,\quad b_n=\CT Q^n.
\]
Matching the three exponents proves these identities directly.

The same exact run found no failures of the following \emph{unproved}
stronger congruences:
\begin{align}
 S(p^2m)&\stackrel{?}{\equiv}p^2S(m)
       \pmod {p^{v_p(m)+3}}\quad(p\text{ odd}),\label{empdouble}\\
 S(pm)&\stackrel{?}{\equiv}pS(m)
       \pmod {p^{v_p(m)+2}}\quad(p\equiv1\pmod4).\label{empsplit}
\end{align}
There were \(989\) tested pairs in \eqref{empdouble}, with \(p^2m\le5000\),
and \(3787\) in \eqref{empsplit}, with \(pm\le5000\).
If \eqref{empdouble} were proved for all its parameters, induction in
steps of two, starting from \eqref{primeprop}, would settle every odd
prime power and hence the entire conjecture. We have not proved it.
The valid transfer \eqref{transferformula}, whose modulus is only \(p^2\),
must not be substituted for \eqref{empdouble}.

\section{Status, provenance, and limits}

The exact unchanged target is currently recorded as
\texttt{OeisA179524.conjecture4}, with \texttt{category research open}
and an unfilled proof, in the inspected formal-conjectures source
\cite{Formal}. The inspected path history first adds the file on
20 August 2026, not 6 November. The historical mathematical source is
already the 2010 conjecture, irrespective of later lists.

Although the OEIS Open paper \cite{OEISOpen} does not mention the
sequence by its identifier, its released benchmark does include it.
The isolated target \texttt{oeis\_a179524\_sun\_conjecture\_2} conjoins
\eqref{target} with the separate prime-square refinement
\[
 S(p)\stackrel{?}{\equiv}
 p\left(7+10\left(\frac{-1}{p}\right)\right)\pmod {p^2}.
\]
The three inspected full-benchmark submissions did not supply accepted
proofs of that conjunction. Omission from the paper's prose is therefore
not evidence of omission from the benchmark.

The full repeated-odd-prime-power problem remains unresolved in this
note. In particular, no general \(p^3\) divisibility is asserted, even
for the inert primes in Theorem~\ref{inert}. The additional small-prime
square checks do not cover arbitrarily large split-prime squares.
The separate inert unweighted branch is proved in
Section~\ref{zeroliftsec}; it must not be identified with the original
weighted target.

Independent derivations checked the digit expansions, high-carry cases,
the CM identification, all-cofactor arguments, and the quadratic
invariant including its exceptional degree-\(p\) coefficient. This is mathematical
review within the research workflow, not external peer review or Lean
formalization. The finite calculations are reproducible in Appendix A.
No claim is made that the partial results, their proof mechanisms, or
their combination are historically first. The known ingredients are
attributed above. The original full-success requirements are not all
met: the full all-\(n\) quantifier is not closed, and nonduplication has
not been established by a comprehensive priority review.

\appendix
\section{Complete exact verification program}

The following standard-library-only Python 3.11-or-later program is contained in this
single \texttt{.tex} source. It computes the original \(a_n\) and \(S(n)\),
not a guessed holonomic recurrence. Its exact base calculations support
the finite-prime corollary; the experimental higher lifts remain
experimental. The recorded output was: no divisibility failures through
\(5000\), no failures in \(989\) double-step tests, and no failures in
\(3787\) split-prime tests.

\begin{lstlisting}
from math import isqrt
from pathlib import Path
import time, json, hashlib, sys
sys.set_int_max_str_digits(0)
N=5000
start=time.monotonic()
a=[]; S=[0]
for n in range(N):
    term=1; an=1
    for k in range(n//2):
        numerator=-4*term*(n-2*k)**2*(n-2*k-1)**2
        denominator=(k+1)**4
        term,rem=divmod(numerator,denominator)
        assert rem==0
        an+=term
    a.append(an)
    S.append(S[-1]+(20*n+17)*an)
    assert S[-1]%(n+1)==0, ('counterexample',n+1)
    if (n+1)%1000==0: print('Exact original sums checked through',n+1,'elapsed',round(time.monotonic()-start,2),flush=True)

def prime(p): return p>=2 and all(p%d for d in range(2,isqrt(p)+1))
def valuation(n,p):
    v=0
    while n%p==0: v+=1; n//=p
    return v

double=[]; split=[]; double_count=split_count=0
for p in range(3,N+1,2):
    if not prime(p): continue
    for m in range(1,N//(p*p)+1):
        mod=p**(valuation(m,p)+3)
        double_count+=1
        if (S[p*p*m]-p*p*S[m])%mod: double.append((p,m,mod))
    if p%4==1:
        for m in range(1,N//p+1):
            mod=p**(valuation(m,p)+2)
            split_count+=1
            if (S[p*m]-p*S[m])%mod: split.append((p,m,mod))

result={'n_max':N,'original_divisibility_failures':[],
        'double_lift_count':double_count,'double_lift_failures':double,
        'split_lift_count':split_count,'split_lift_failures':split,
        'elapsed_seconds':round(time.monotonic()-start,3),
        'a_sha256':hashlib.sha256(('\n'.join(map(str,a))+'\n').encode()).hexdigest()}
print(json.dumps(result,indent=2),flush=True)
Path(__file__).with_name('exact_lift_test_results.json').write_text(json.dumps(result,indent=2)+'\n')
\end{lstlisting}

\begin{thebibliography}{99}
\raggedright
\bibitem{Sun2010}
Z.-W. Sun, \emph{On Ap\'ery numbers and generalized central trinomial
coefficients}, arXiv:1006.2776v11, 19 August 2010.
Definition of \(W_n(b,c)\), p.~29; Conjecture 5.8(i)--(ii), pp.~32--33.
\url{https://arxiv.org/abs/1006.2776v11}.

\bibitem{OEIS}
Z.-W. Sun, \emph{A179524}, The On-Line Encyclopedia of Integer Sequences,
entry dated 17 July 2010, conjecture attributed to 1 July 2010;
accessed 7 November 2026.
\url{https://oeis.org/A179524}.

\bibitem{SunRefinement}
Z.-W. Sun, \emph{A refinement of a congruence result by van Hamme and
Mortenson}, Illinois Journal of Mathematics \textbf{56} (2012), no.~3,
967--979. Theorem 1.2.
\url{https://arxiv.org/abs/1011.1902v6}.

\bibitem{Duke}
W. Duke, \emph{Modular functions and the uniform distribution of CM
points}, Mathematische Annalen \textbf{334} (2006), 241--252.
The first page gives \(j(i\sqrt5)=632000+282880\sqrt5\).
\url{https://math.ucla.edu/~wdduke/preprints/modud.pdf}.

\bibitem{EOY}
N. Elkies, K. Ono and T. Yang, \emph{Reduction of CM elliptic curves and
modular function congruences}, International Mathematics Research
Notices (2005), no.~44, 2695--2707.
The introduction recalls Deuring's nonsplit-prime reduction theorem.
\url{https://people.math.wisc.edu/~tonghaiyang/deuring.pdf}.

\bibitem{Silverman}
J. H. Silverman, \emph{The Arithmetic of Elliptic Curves}, second edition,
Graduate Texts in Mathematics 106, Springer, 2009.
Theorem V.4.1(a), p.~148.

\newpage
\bibitem{Formal}
Google DeepMind and contributors, \emph{formal-conjectures},
\texttt{FormalConjectures/OEIS/179524.lean}, inspected 7 November 2026.
Initial path commit \texttt{1bb238f2be65a7d796fb29a15f1c89f764062515},
20 August 2026.
\url{https://github.com/google-deepmind/formal-conjectures/blob/main/FormalConjectures/OEIS/179524.lean}.

\bibitem{OEISOpen}
T. Adamczewski, \emph{OEIS Open: How many conjectures can language
models turn into theorems?}, arXiv:2608.11941v2, 13 August 2026.
\url{https://arxiv.org/abs/2608.11941v2}.
Released exact target:
\url{https://github.com/epoch-research/LeanOpenProblems/blob/main/apn/data/oeis/Isolated/oeis_a179524_sun_conjecture_2.lean}.

\bibitem{SunApery}
Z.-W. Sun, \emph{On sums of Ap\'ery polynomials and related congruences},
Journal of Number Theory \textbf{132} (2012), 2673--2699.
Theorem 1.1 and Remark 1.1.
\url{https://arxiv.org/abs/1101.1946v4}.

\bibitem{GuoZeng}
V. J. W. Guo and J. Zeng, \emph{Proof of some conjectures of Z.-W. Sun
on congruences for Ap\'ery polynomials}, Journal of Number Theory
\textbf{132} (2012), 1731--1740; arXiv:1101.0983.
Theorem 1.3.
\url{https://arxiv.org/abs/1101.0983}.

\bibitem{Babanskyy}
O. Babanskyy, \emph{Sun's supercongruence for two Ap\'ery-like binomial
sums at discriminant \(-20\)}, public manuscript, 2026.
Section 2, Lemmas 2.1--2.2, pp.~2--3.
Pinned repository commit
\texttt{4caad01fa469c86a1a95b930293354d843b79d64};
the inspected repository history dates from 3 October 2026.
\url{https://github.com/aconsciousfractal/Sun-3-7-Supercongruence/blob/4caad01fa469c86a1a95b930293354d843b79d64/paper/Sun_3_7_Supercongruence.pdf}.

\end{thebibliography}
\end{document}

